\begin{enumerate}

\item Se obtiene beta de raíz con componente tangencial -U.

\[\fitmath{\mathrm{bh}=\arctan\left(\frac{-\Omega\cdot \mathrm{rh}}{\mathrm{cx}}\right)}\]

\[\fitmath{\mathrm{bh}=\arctan\left(\frac{-(1128.9)\cdot (0.198)}{(157)}\right)\approx -0.9584674\,\mathrm{rad}}\]

\item Se obtiene beta de punta.

\[\fitmath{\mathrm{bt}=\arctan\left(\frac{-\Omega\cdot \mathrm{rt}}{\mathrm{cx}}\right)}\]

\[\fitmath{\mathrm{bt}=\arctan\left(\frac{-(1128.9)\cdot (0.314)}{(157)}\right)\approx -1.153855\,\mathrm{rad}}\]

\item La torsión es una diferencia algebraica de orientaciones.

\[\fitmath{\mathrm{twist}=\arctan\left(\frac{-\Omega\cdot \mathrm{rt}}{\mathrm{cx}}\right)-\arctan\left(\frac{-\Omega\cdot \mathrm{rh}}{\mathrm{cx}}\right)}\]

\[\fitmath{\mathrm{twist}=\arctan\left(\frac{-(1128.9)\cdot (0.314)}{(157)}\right)-\arctan\left(\frac{-(1128.9)\cdot (0.198)}{(157)}\right)\approx -0.1953875\,\mathrm{rad}}\]

\end{enumerate}
